\documentclass[12pt,a4paper]{article}
\usepackage{graphicx}
\usepackage{geometry}
\usepackage{amsmath}
\usepackage{fancyhdr}
\usepackage{setspace}

\geometry{margin=1in}
\onehalfspacing

\pagestyle{fancy}
\fancyhf{}
\rhead{COMETFWC054}
\lhead{Digital Logic}
\cfoot{\thepage}

\begin{document}

\begin{center}

\includegraphics[width=4cm]{assets/logo.jpg}

\vspace{0.5cm}

{\Large \textbf{ARM 2010 EC Question 41}}

\vspace{0.5cm}

\textbf{Name: KESHAVA REDDY V} \\
\textbf{ID: COMETFWC054}
\end{center}

\vspace{1cm}

\section*{Question}

For any set of inputs, $A$ and $B$, the following circuits give the same output $Q$, except one. Which one is it? (GATE PH 2010)

\begin{center}
\includegraphics[width=0.8\textwidth]{pics/q41.jpg}
\end{center}

\section*{Question Analysis}

We analyze each circuit by deriving its Boolean expression.

\subsection*{Option (a)}

The circuit represents:

\[
Q = (A + B) \cdot \overline{B}
\]

Using Boolean algebra:

\[
Q = A\overline{B} + B\overline{B}
\]

Since $B\overline{B} = 0$,

\[
Q = A\overline{B}
\]

\subsection*{Option (b)}

This circuit directly represents:

\[
Q = A \cdot \overline{B}
\]

\subsection*{Option (c)}

After simplifying the logic network:

\[
Q = A\overline{B}
\]

\subsection*{Option (d)}

This circuit represents:

\[
Q = \overline{A} + B
\]

This expression is different from $A\overline{B}$.

\section*{Conclusion from Analysis}

Options (a), (b), and (c) produce:

\[
Q = A\overline{B}
\]

Option (d) produces:

\[
Q = \overline{A} + B
\]

Hence, the incorrect circuit is:

\[
\boxed{\text{Option (d)}}
\]

\section*{Truth Table}

\begin{center}
\begin{tabular}{|c|c|c|c|}
\hline
A & B & $A\overline{B}$ & $\overline{A}+B$ \\
\hline
0 & 0 & 0 & 1 \\
0 & 1 & 0 & 1 \\
1 & 0 & 1 & 0 \\
1 & 1 & 0 & 1 \\
\hline
\end{tabular}
\end{center}

The table clearly shows Option (d) differs from the other circuits.

\section*{Hardware Implementation}

The Boolean function $Q = A\overline{B}$ is implemented using Raspberry Pi Pico 2W.

The microcontroller generates all input combinations using a loop. For each input:

\[
Q = A\overline{B}
\]

If $Q = 1$, the LED turns ON and the seven-segment display shows 1.  
If $Q = 0$, the LED turns OFF and the display shows 0.

\section*{Required Components}

\begin{itemize}
\item Raspberry Pi Pico 2W
\item Seven Segment Display
\item LED
\item 330$\Omega$ resistors
\item Breadboard
\item Jumper wires
\end{itemize}

\section*{Experimental Setup}

\begin{center}
\includegraphics[width=0.7\textwidth]{pics/setup1.jpg}

\vspace{0.5cm}

\includegraphics[width=0.7\textwidth]{pics/setup2.jpg}
\end{center}

\section*{Final Conclusion}

The circuits in options (a), (b), and (c) are equivalent and implement:

\[
Q = A\overline{B}
\]

Option (d) implements a different function:

\[
Q = \overline{A} + B
\]

Therefore, the correct answer is:

\[
\boxed{\text{Option (d)}}
\]

\end{document}
